% Generated from validated Article Draft Result. Do not hand-edit; revise the structured draft instead.
\section{回答から検索・操作・接続へ}
\label{pkg:action-connectivity}
\noindent\textbf{2024年後半は、生成AIの外部世界との接続面が急に増えた。検索、GUI操作、structured tool interface、protocol、agent evaluationを一つの『Agent』に押し込めず、異なるexecution primitiveとして分解する。}\autocite{src-6cafd8cf5668,src-9e4ef9c155b6}

SearchGPTのprototypeとChatGPT searchの製品化は、同じ検索機能の一枚岩のreleaseではない。検索を回答生成へ接続するsurfaceが、実験的な入口から日常的なproduct capabilityへ移る過程として、event identityを分けて追う必要がある。\autocite{src-6cafd8cf5668,src-bd5105ccfa25}


Claude 3.5 SonnetのComputer Useは、modelが画面上の操作へ踏み込むpublic betaという別種のaction surfaceを示した。一方、structured tool interfaceやGemini系のtool use、OpenAI側のRealtime/API周辺更新は、GUI操作とは異なる機械可読な実行経路を広げた。どれも外部作用を可能にするが、入力形式、権限、失敗モード、評価方法は同じではない。\autocite{src-9e4ef9c155b6,src-a38408fd5407,src-e283b9c34207,src-ed98251ab05d}


MCPはmodel capabilityそのものではなく、contextやtoolsを接続するprotocol layerとして位置付けるべきである。同時にAppWorldやToolSandbox、OS操作・tool-use系の研究は、接続面が増えるほど『できるか』だけでなく『長いtaskをどの程度再現可能に実行できるか』を測る評価基盤が必要になることを示した。protocolの登場とagent reliabilityの解決は別の問題である。\autocite{src-0d67af76cfef,src-395ebe20f9b0,src-5ad3266ea448,src-7afacff558f5}


\begin{claimboundary}
Computer Use、tool-use benchmark、MCP、searchを同じ『agent完成度』の指標に変換しない。一次資料はそれぞれのrelease、interface、評価設定を確定するが、異なるsurface間の性能優位やproduction reliabilityを独立再現したものではない。\autocite{src-9e4ef9c155b6}
\end{claimboundary}
